\chapter[Appendix A]{Appendix A}
\label{sec:appendix}

\section{Device Specifications}

\subsection{Host system}

\begin{verbatim}
hostnamectl
Static hostname: anji-Nitro-AN515-58
Icon name: computer-laptop
Chassis: laptop
Machine ID: 07bcf4e12b18401998b1ced41087efdf
Boot ID: df0e8cb23e204d1080c42530bfeb49ad
Operating System: Ubuntu 26.04 LTS

Kernel: Linux 7.0.0-28-generic
Architecture: x86-64
Hardware Vendor: Acer
Hardware Model: Nitro AN515-58
Hardware SKU: 0000000000000000
Hardware Version: V2.21
Firmware Version: V2.21
Firmware Date: Tue 2026-06-09
Firmware Age: 1month 4w
\end{verbatim}

\subsection{Test Drive 1: NVMe SSD}
\begin{verbatim}
sudo smartctl -a /dev/sda

smartctl 7.5 2025-04-30 r5714[x86_64-linux-7.0.0-28-generic](local build)
Copyright(C)2002-25, Bruce Allen, Christian Franke,www.smartmontools.org

=== START OF INFORMATION SECTION ===
Model Number:                       PSELN128GA87MC0
Serial Number:                      0524115009487
Firmware Version:                   HPE4BP01
PCI Vendor/Subsystem ID:            0x1987
IEEE OUI Identifier:                0x6479a7
Controller ID:                      1
NVMe Version:                       1.4
Number of Namespaces:               1
Namespace 1 Size/Capacity:          128,035,676,160 [128 GB]
Namespace 1 Formatted LBA Size:     512
Namespace 1 IEEE EUI-64:            6479a7 8d8f00250f
Local Time is:                      Thu Aug  6 12:04:21 2026 CEST
Firmware Updates (0x04):            2 Slots
Optional Admin Commands (0x0017):   Security Format Frmw_DL Self_Test
Optional NVM Commands (0x0056):Wr_Unc DS_Mngmt Sav/Sel_Feat Timestmp
Log Page Attributes (0x1e):Cmd_Eff_Lg Ext_Get_Lg Telmtry_Lg Pers_Ev_Lg
Maximum Data Transfer Size:64 Pages
Warning  Comp. Temp. Threshold:     83 Celsius
Critical Comp. Temp. Threshold:     85 Celsius

Supported Power States
St Op     Max   Active     Idle   RL RT WL WT  Ent_Lat  Ex_Lat
0 +     4.30W       -        -    0  0  0  0        0       0
1 +     3.00W       -        -    1  1  1  1        0       0
2 +     1.50W       -        -    2  2  2  2        0       0
3 -   0.1000W       -        -    3  3  3  3     5000   10000
4 -   0.0025W       -        -    4  4  4  4     5000   45000

Supported LBA Sizes (NSID 0x1)
Id Fmt  Data  Metadt  Rel_Perf
0 +     512       0         1
1 -    4096       0         0

=== START OF SMART DATA SECTION ===
SMART overall-health self-assessment test result: PASSED

SMART/Health Information (NVMe Log 0x02, NSID 0xffffffff)
Critical Warning:                   0x00
Temperature:                        57 Celsius
Available Spare:                    100%
Available Spare Threshold:          5%
Percentage Used:                    0%
Data Units Read:                    1,417,462 [725 GB]
Data Units Written:                 393,172 [201 GB]
Host Read Commands:                 11,854,846
Host Write Commands:                2,109,150
Controller Busy Time:               36
Power Cycles:                       72
Power On Hours:                     11
Unsafe Shutdowns:                   24
Media and Data Integrity Errors:    0
Error Information Log Entries:      64
Warning  Comp. Temperature Time:    0
Critical Comp. Temperature Time:    0
Temperature Sensor 1:               57 Celsius
\end{verbatim}

\subsection{Test Drive 2: SATA SSD}
\begin{verbatim}
sudo hdparm -I /dev/sdb

/dev/sdb:

ATA device, with non-removable media
Model Number:       EDILOCA  ES106 128GB

Serial Number:      26060047623000204

Firmware Revision:  HC2705H3
Media Serial Num:

Media Manufacturer:
Transport:Serial, ATA8-AST,SATA 1.0a,SATA II Extensions,SATA Rev 2.5,SATA Rev 2.6,SATA Rev 3.0
Standards:
Used: unknown (minor revision code 0x0110)
Supported: 9 8 7 6 5
Likely used: 9
Configuration:
Logical		max	current
cylinders	16383	16383
heads		16	16
sectors/track	63	63
--
CHS current addressable sectors:    16514064
LBA    user addressable sectors:   250069680
LBA48  user addressable sectors:   250069680
Logical  Sector size:                   512 bytes
Physical Sector size:                   512 bytes
Logical Sector-0 offset:                  0 bytes
device size with M = 1024*1024:      122104 MBytes
device size with M = 1000*1000:      128035 MBytes (128 GB)
cache/buffer size  = unknown
Form Factor: 2.5 inch
Nominal Media Rotation Rate: Solid State Device
Capabilities:
LBA, IORDY(can be disabled)
Queue depth: 32
Standby timer values: spec'd by Standard, no device specific minimum
R/W multiple sector transfer: Max = 8 Current = 1
DMA: mdma0 mdma1 mdma2 udma0 udma1 udma2 udma3 udma4 udma5 *udma6
Cycle time: min=120ns recommended=120ns
PIO: pio0 pio1 pio2 pio3 pio4
Cycle time: no flow control=120ns  IORDY flow control=120ns
Commands/features:
Enabled	Supported:
*	SMART feature set
Security Mode feature set
*	Power Management feature set
*	Write cache
*	Look-ahead
*	WRITE_BUFFER command
*	READ_BUFFER command
*	NOP cmd
*	48-bit Address feature set
*	Mandatory FLUSH_CACHE
*	FLUSH_CACHE_EXT
*	SMART error logging
*	SMART self-test
*	General Purpose Logging feature set
*	WRITE_{DMA|MULTIPLE}_FUA_EXT
*	64-bit World wide name
*	{READ,WRITE}_DMA_EXT_GPL commands
*	Gen1 signaling speed (1.5Gb/s)
*	Gen2 signaling speed (3.0Gb/s)
*	Gen3 signaling speed (6.0Gb/s)
*	Native Command Queueing (NCQ)
*	Phy event counters
*	Host automatic Partial to Slumber transitions
*	Device automatic Partial to Slumber transitions
*	READ_LOG_DMA_EXT equivalent to READ_LOG_EXT
DMA Setup Auto-Activate optimization
*	Software settings preservation
*	SANITIZE feature set
*	BLOCK_ERASE_EXT command
*	WRITE BUFFER DMA command
*	READ BUFFER DMA command
*	Data Set Management TRIM supported (limit 8 blocks)
Security:
Master password revision code = 65534
supported
not	enabled
not	locked
not	frozen
not	expired: security count
supported: enhanced erase
2min for SECURITY ERASE UNIT. 2min for ENHANCED SECURITY ERASE UNIT.
Logical Unit WWN Device Identifier:0000000000000000
NAA		: 0
IEEE OUI	: 000000
Unique ID	: 000000000
Checksum: correct
\end{verbatim}


\section{Experiment 1}

\begin{verbatim}

sudo mkdir -p /mnt/test/test_data
sudo openssl rand -out "/mnt/test/test_data/
inode_test001.dat" 104857600


sudo find "/mnt/test/test_data" -type f -exec stat -c '%n 
| Inode: %i' {} \;
/mnt/test/test_data/inode_test001.dat | Inode: 6291458

sudo dumpe2fs -h /dev/sda1 | grep -E "Journal|Inode"

sudo rm -rf "/mnt/test/test_data"/*

sudo debugfs -R "stat <6291458>" /dev/sda1

sudo dumpe2fs -h /dev/sda1 | grep -E "Journal|Inode"

sudo dumpe2fs -h /dev/sda1 | grep -E "Journal|Inode"

sudo ext4magic /dev/sda1 -j /dev/sda1 -I 6291458 -d /tmp/recovery/

sudo ext4magic /dev/sda1 -j /dev/sda1 -I 6291458 -B 620 -d 
/tmp/recovery/

sudo debugfs /dev/sda1
debugfs: ncheck 6291458

sudo debugfs -R "imap <6291458>" /dev/sda1

sudo debugfs -R "logdump" /dev/sda1
sudo debugfs -R "logdump -a" /dev/sda1

sudo dd if=/dev/sda1 bs=4096 skip=33388 count=1 2>/dev/null | xxd 
| head -20
sudo dd if=/dev/sda1 bs=4096 skip=33389 count=1 2>/dev/null | xxd 
| head -20

sudo dd if=/dev/sda1 of=~/inode_analysis/journal_tx61.bin bs=4096 
skip=15237743 count=1
sudo dd if=~/inode_analysis/journal_tx61.bin 
of=~/inode_analysis/inode_tx61.bin bs=1 skip=256 count=256
\end{verbatim}

\section{Experiment 2}
\subsection{Experimental Setup}
\begin{verbatim}
 # 1. Unmount if mounted
sudo umount /dev/sda1 2>/dev/null

# 2. Wipe the drive (quick wipe)
sudo dd if=/dev/zero of=/dev/sda bs=1M count=100 status=progress

# 3. Create GPT partition table
sudo parted /dev/sda mklabel gpt

# 4. Create single ext4 partition
sudo parted /dev/sda mkpart primary ext4 0% 100%

# 5. Reload partition table
sudo partprobe /dev/sda
sleep 2

# 6. Format as ext4
sudo mkfs.ext4 -F /dev/sda1

# 7. Create mount point
sudo mkdir -p /mnt/{MountPoint}

\end{verbatim}

\subsection{Trim-on and Trim-off Conditions}
\begin{verbatim}

# TRIM ON
sudo mount -o discard /dev/sda1 /mnt/{MountPoint}

# TRIM OFF
sudo mount -o nodiscard /dev/sda1 /mnt/{MountPoint}

\end{verbatim}

\subsection{Data Generation Script}

\begin{lstlisting}[language=bash, caption={Data Generation Script}]

#!/bin/bash
# Data generation script

MOUNT_POINT=$1
DATASET_NUM=$2
UNIQUE_ID=$(openssl rand -hex 8)

FOLDER_NAME="TEST_DATA_${DATASET_NUM}_${UNIQUE_ID}"
mkdir -p "${MOUNT_POINT}/${FOLDER_NAME}"

# File 1: Small text file (51 signatures)
cat > "${MOUNT_POINT}/${FOLDER_NAME}/small.txt" << ENDTEXT
UNIQUE_SIGNATURE_DS${DATASET_NUM}_SMALL
$(for i in {1..50}; do echo "UNIQUE_SIGNATURE_DS${DATASET_NUM}_SMALL"; done)
ENDTEXT

# File 2: Medium text file (5,001 signatures)
cat > "${MOUNT_POINT}/${FOLDER_NAME}/medium.txt" << ENDTEXT
UNIQUE_SIGNATURE_DS${DATASET_NUM}_MEDIUM
$(for i in {1..5000}; do echo "UNIQUE_SIGNATURE_DS${DATASET_NUM}_MEDIUM"; done)
ENDTEXT

# File 3: Large binary file (10MB + 1 signature)
dd if=/dev/urandom of="${MOUNT_POINT}/${FOLDER_NAME}/large.bin" bs=1M count=10 status=none
echo "UNIQUE_SIGNATURE_DS${DATASET_NUM}_LARGE" >> "${MOUNT_POINT}/${FOLDER_NAME}/large.bin"

# File 4: 500 small text files (500 signatures)
for i in {1..500}; do
    echo "UNIQUE_SIGNATURE_DS${DATASET_NUM}_SMALL_${i}" > "${MOUNT_POINT}/${FOLDER_NAME}/tiny_${i}.txt"
done

# File 5: Binary image file (no signatures)
dd if=/dev/urandom of="${MOUNT_POINT}/${FOLDER_NAME}/image.jpg" bs=1K count=50 status=none

# File 6: Binary PDF file (no signatures)
dd if=/dev/urandom of="${MOUNT_POINT}/${FOLDER_NAME}/document.pdf" bs=1K count=30 status=none

# Create checksum file for verification
md5sum "${MOUNT_POINT}/${FOLDER_NAME}"/* > "${MOUNT_POINT}/${FOLDER_NAME}/checksums_${DATASET_NUM}.md5"

echo "Created dataset $DATASET_NUM in ${MOUNT_POINT}/${FOLDER_NAME}"

\end{lstlisting}

\subsection{Command List}
\subsubsection*{Phase A: TRIM ON Procedure}
\begin{enumerate}
    
\item{Mount with TRIM enabled:}

\begin{verbatim}
sudo mount -o discard /dev/sda1 /mnt/ssd_test
\end{verbatim}

\item{Create Datasets 1-5 with unique signatures:}

\begin{verbatim}
for i in {1..5}; do
    ./create_dataset.sh /mnt/ssd_test $i
done
\end{verbatim}

\item{Delete Datasets 2 and 4:}

\begin{verbatim}
sudo rm -rf /mnt/ssd_test/TEST_DATA_2_*
sudo rm -rf /mnt/ssd_test/TEST_DATA_4_*
\end{verbatim}

\item{Force TRIM:}

\begin{verbatim}
sudo fstrim -v /mnt/ssd_test
\end{verbatim}

\item{Create Overwrite Datasets 6-8:}

\begin{verbatim}
for i in {6..8}; do
    ./create_dataset.sh /mnt/ssd_test $i
done
\end{verbatim}

\item{Unmount the drive:}

\begin{verbatim}
sudo umount /mnt/ssd_test
\end{verbatim}

\end{enumerate}

\subsubsection*{Phase B: TRIM OFF Procedure}

\begin{enumerate}
    
\item{Mount with TRIM enabled:}

\begin{verbatim}
sudo mount -o discard /dev/sda1 /mnt/ssd_test
\end{verbatim}

\item{Create Datasets 1-5 with unique signatures:}

\begin{verbatim}
for i in {1..5}; do
    ./create_dataset.sh /mnt/ssd_test $i
done
\end{verbatim}

\item{Delete Datasets 2 and 4:}

\begin{verbatim}
sudo rm -rf /mnt/ssd_test/TEST_DATA_2_*
sudo rm -rf /mnt/ssd_test/TEST_DATA_4_*
\end{verbatim}

\item{Force TRIM:}

\begin{verbatim}
sudo fstrim -v /mnt/ssd_test
\end{verbatim}

\item{Create Overwrite Datasets 6-8:}

\begin{verbatim}
for i in {6..8}; do
    ./create_dataset.sh /mnt/ssd_test $i
done
\end{verbatim}

\item{Unmount the drive:}

\begin{verbatim}
sudo umount /mnt/ssd_test
\end{verbatim}

\end{enumerate}

\subsubsection*{Phase B: TRIM OFF Procedure}
\begin{enumerate}

\item{Mount without TRIM:}

\begin{verbatim}
sudo mount -o nodiscard /dev/sda1 /mnt/ssd_test
\end{verbatim}

\item{Create Datasets 9-13:}

\begin{verbatim}
for i in {9..13}; do
    ./create_dataset.sh /mnt/ssd_test $i
done
\end{verbatim}

\item{Delete Datasets 10 and 12:}

\begin{verbatim}
sudo rm -rf /mnt/ssd_test/TEST_DATA_10_*
sudo rm -rf /mnt/ssd_test/TEST_DATA_12_*
\end{verbatim}

\item{NO TRIM - deliberately skipped}

\item{Create Overwrite Datasets 14-16:}

\begin{verbatim}
for i in {14..16}; do
    ./create_dataset.sh /mnt/ssd_test $i
done
\end{verbatim}

\item{Unmount the drive:}

\begin{verbatim}
sudo umount /mnt/ssd_test
\end{verbatim}

\end{enumerate}

\begin{verbatim}
for i in {1..16}; do
count=$(grep -r "UNIQUE_SIGNATURE_DS${i}_". 2>/dev/null | wc -l)
echo "Dataset $i: $count files found"
done
\end{verbatim}

Limitations and workarounds:
\begin{enumerate}
    \item Identify the USB enclosure's chip and obtain its vendor and product ID using the command:
        \begin{verbatim}
        lsusb
        \end{verbatim}

    \item Created udev rule to enable unmap
        \begin{verbatim}
        sudo nano /etc/udev/rules.d/90-usb-discard.rules
        echo 'ACTION=="add|change", ATTRS{idVendor}=="152d", 
        ATTRS{idProduct}=="a583", SUBSYSTEM=="scsi_disk", 
        ATTR{provisioning_mode}="unmap"' | sudo tee 
        /etc/udev/rules.d/90-usb-discard.rules
        \end{verbatim}

    \item Reloaded udev rules and replugged the device
        \begin{verbatim}
        sudo udevadm control --reload-rules
        sudo udevadm trigger
        \end{verbatim}

    \item Verified TRIM was recognized. \texttt{DISC-GRAN and DISC-MAX} values should be non-zero. 
        \begin{verbatim}
        lsblk --discard
        # DISC-GRAN=4K, DISC-MAX=4G (TRIM recognized)
        \end{verbatim}

\end{enumerate}

\section{Experiment 3: Full Script}

\begin{lstlisting}[language=bash, caption={PHASE 1: wipe drive and prepare for experiment}]

# ============================================
# PHASE 1: wipe drive
# ============================================

TEST_DEVICE="/dev/sda"
MOUNT_POINT="/mnt/test_ssd"

echo ""
echo "=== PHASE 1: Wipe Drive ==="
echo "WARNING: This will erase ALL data on $TEST_DEVICE"
read -p "Type 'y' to continue: " confirm

if [ "$confirm" != "y" ]; then
    echo "Aborted."
    exit 1
fi


sudo umount $MOUNT_POINT 2>/dev/null
echo "Overwriting $TEST_DEVICE with zeros..."
sudo dd if=/dev/zero of=$TEST_DEVICE bs=1M count=100 status=progress

echo "Creating partition table..."
sudo parted -s $TEST_DEVICE mklabel gpt
sudo parted -s $TEST_DEVICE mkpart primary ext4 0% 100%
sudo mkfs.ext4 -F ${TEST_DEVICE}1

sudo mkdir -p $MOUNT_POINT
sudo mount ${TEST_DEVICE}1 $MOUNT_POINT
sudo chmod 777 $MOUNT_POINT
echo "Wipe complete"
echo "Trigger udev rules to enable trim"
sudo udevadm control --reload-rules
sudo udevadm trigger
\end{lstlisting}

\begin{lstlisting}[language=bash, caption={PHASE 2: Configure andFill Drive to Target Utilisation Level}]
# ============================================
# SIMPLE CONFIGURATION
# ============================================

LOG_DIR="$HOME/trim_experiment/logs_nvme"
TEST_DIR="$MOUNT_POINT/test_data"
FILLER_DIR="$MOUNT_POINT/filler_data"

echo "Input Fill Level:"
read -p "Fill level (%): " FILL_LEVEL_INPUT

# test parameters
FILL_LEVEL=$FILL_LEVEL_INPUT                
RECOVERY_WAITS=(0 5 15 30)            
ITERATIONS=${#RECOVERY_WAITS[@]}

echo ""
echo "Configuration:"
echo "  Fill level: ${FILL_LEVEL}%"
echo "  Test data: ${TEST_DATA_SIZE_MB}MB"
echo "  Recovery waits: ${RECOVERY_WAITS[@]} minutes"
echo "  Total iterations: ${ITERATIONS}"
echo ""

mkdir -p $LOG_DIR
sudo mkdir -p $TEST_DIR $FILLER_DIR 


# ============================================
# PHASE 2: Fill Drive to FILL_LEVEL%
# ============================================

echo ""
echo "=== PHASE 2: Fill Drive to ${FILL_LEVEL}% ==="

# Calculate how much to fill
TOTAL_KB=$(df --output=size "$MOUNT_POINT" | tail -1)
USED_KB=$(df --output=used "$MOUNT_POINT" | tail -1)
TARGET_KB=$((TOTAL_KB * FILL_LEVEL / 100))
FILL_KB=$((TARGET_KB - USED_KB))

echo "Need to write $((FILL_KB / 1024)) MB to reach ${FILL_LEVEL}%"

# Create filler files
written_kb=0
file_count=0
while [ $written_kb -lt $FILL_KB ]; do
    file_name="$FILLER_DIR/filler_$(printf "%04d" $file_count).dat"
    this_size=$((FILL_KB - written_kb))
    if [ $this_size -gt $((1024 * 100)) ]; then  # 100MB chunks
        this_size=$((1024 * 100))
    fi
    
    sudo dd if=/dev/urandom of="$file_name" bs=1K count=$this_size status=none
    written_kb=$((written_kb + this_size))
    file_count=$((file_count + 1))
    
    echo -ne "\r  Progress: $((written_kb * 100 / FILL_KB))%"
done
echo ""
echo "Created $file_count filler files"

# Check actual fill level
NEW_USED=$(df --output=used "$MOUNT_POINT" | tail -1)
NEW_PERCENT=$((NEW_USED * 100 / TOTAL_KB))
echo "Actual fill level: ${NEW_PERCENT}%"
\end{lstlisting}

\begin{lstlisting}[language=bash, caption={PHASE 3-6: Main Experiment Loop}]
# ============================================
# MAIN EXPERIMENT LOOP - Each iteration tests complete cycle
# ============================================

echo ""
echo "=== PHASE 3-6: MAIN EXPERIMENT ($ITERATIONS iterations) ==="
echo "Each iteration: Create -> Delete -> TRIM -> Recover"
echo ""

# Initialize results files
if [ ! -f "$LOG_DIR/experiment_results.csv" ]; then
    echo "run_timestamp,iteration,wait_minutes,fill_level,real_time,user_time,sys_time,trimmed_bytes,recovered_files" > $LOG_DIR/experiment_results.csv
fi

# Main loop - one iteration per recovery wait time
for ITER in $(seq 1 $ITERATIONS); do
    # Get the wait time for THIS iteration (0-indexed array)
    WAIT_INDEX=$((ITER - 1))
    CURRENT_WAIT=${RECOVERY_WAITS[$WAIT_INDEX]}
    
    echo "=========================================="
    echo "ITERATION $ITER of $ITERATIONS - Wait: ${CURRENT_WAIT} minutes"
    echo "=========================================="
    
# ============================================
# PHASE 3: Create Test Data
# ============================================

echo ""
echo "=== [ITER $ITER - PHASE 3]: Create Test Data ==="

# Create UNIQUE test data for THIS iteration
UNIQUE_ID=$(openssl rand -hex 8)
ITER_TEST_DIR="$TEST_DIR/TEST_DATA_${ITER}_${UNIQUE_ID}"
sudo mkdir -p "$ITER_TEST_DIR"

echo "Creating dataset in: $ITER_TEST_DIR"
echo ""

# ============================================
# FILE 1: Small text file
# ============================================

echo "Creating small text file..."

{
    echo "UNIQUE_SIGNATURE_DS${ITER}_SMALL"
    for i in {1..50}; do
        echo "UNIQUE_SIGNATURE_DS${ITER}_SMALL"
    done
} | sudo tee "$ITER_TEST_DIR/small.txt" > /dev/null

echo "small.txt created"


# ============================================
# FILE 2: Medium text file
# ============================================

echo "Creating medium text file..."

{
    echo "UNIQUE_SIGNATURE_DS${ITER}_MEDIUM"
    for i in {1..5000}; do
        echo "UNIQUE_SIGNATURE_DS${ITER}_MEDIUM"
    done
} | sudo tee "$ITER_TEST_DIR/medium.txt" > /dev/null
echo "medium.txt created"


# ============================================
# FILES 3-7: LARGE VALID ZIP FILES
# Approximately 1 GiB each
# ============================================

echo ""
echo "Creating large ZIP dataset (~2.5 GiB each)..."


    TEMP_FILE="$ITER_TEST_DIR/.temporary_data.bin"
    ZIP_FILE="$ITER_TEST_DIR/large_archive.zip"

    sudo dd if=/dev/urandom \
        of="$TEMP_FILE" \
        bs=1M \
        count=2560 \
        status=progress

  
    echo "UNIQUE_SIGNATURE_DS${ITER}_LARGE" | \
        sudo tee -a "$TEMP_FILE" > /dev/null

    sudo zip -0 -q "$ZIP_FILE" "$TEMP_FILE"

    # Remove the temporary uncompressed file.
    sudo rm -f "$TEMP_FILE"


echo ""
echo "Created large ZIP archive"


# ============================================
# FILE 8: MANY SMALL TEXT FILES
# ============================================

echo ""
echo "Creating 500 small text files..."

for i in {1..500}; do

    sudo bash -c "echo 'UNIQUE_SIGNATURE_DS${ITER}_SMALL_${i}' > '$ITER_TEST_DIR/tiny_${i}.txt'"

done

echo "Created 500 tiny text files"


# ============================================
# FILE 9: VALID JPEG IMAGE
# ============================================

echo ""
echo "Creating JPEG file..."

# Minimal valid 1x1 JPEG image.
JPEG_BASE64="/9j/4AAQSkZJRgABAQAAAQABAAD/2wBDAP/////////////////////////////////
/////////////////////////////////////////////////////2wBDAf/////////////////////
/////////////////////////////////////////////////////////////////wAARCAABAAEDASI
AAhEBAxEB/8QAFQABAQAAAAAAAAAAAAAAAAAAAAX/xAAUEAEAAAAAAAAAAAAAAAAAAAAA/9oADAMBAAI
QAxAAAAH/AP/EABQQAQAAAAAAAAAAAAAAAAAAACD/2gAIAQEAAT8Af//EABQRAQAAAAAAAAAAAAAAAAA
AABD/2gAIAQIBAT8Af//EABQRAQAAAAAAAAAAAAAAAAAAABD/2gAIAQMBAT8Af//Z"

echo "$JPEG_BASE64" | base64 -d | sudo tee "$ITER_TEST_DIR/image.jpg" > /dev/null

# Add a dataset-specific text signature after the JPEG.
# This does not affect the JPEG header and gives us a searchable
# experiment-specific marker.
echo "UNIQUE_SIGNATURE_DS${ITER}_IMAGE" | \
    sudo tee -a "$ITER_TEST_DIR/image.jpg" > /dev/null

echo "image.jpg created"


# ============================================
# FILE 10: VALID PDF
# ============================================

echo ""
echo "Creating PDF file..."

sudo bash -c "cat > '$ITER_TEST_DIR/document.pdf' << ENDPDF
%PDF-1.4
1 0 obj
<< /Type /Catalog /Pages 2 0 R >>
endobj
2 0 obj
<< /Type /Pages /Kids [3 0 R] /Count 1 >>
endobj
3 0 obj
<< /Type /Page /Parent 2 0 R /MediaBox [0 0 612 792] /Contents 4 0 R >>
endobj
4 0 obj
<< /Length 120 >>
stream
BT
/F1 24 Tf
100 700 Td
(UNIQUE_SIGNATURE_DS${ITER}_PDF) Tj
0 -40 Td
(Experimental secure deletion dataset) Tj
ET
endstream
endobj
xref
0 5
0000000000 65535 f
0000000009 00000 n
0000000058 00000 n
0000000115 00000 n
0000000210 00000 n
trailer
<< /Size 5 /Root 1 0 R >>
startxref
380
%%EOF
ENDPDF"

echo "document.pdf created"

# ============================================
# FILE 11: CHECKSUMS
# ============================================

echo ""
echo "Computing MD5 checksums..."

sudo bash -c "find '$ITER_TEST_DIR' -maxdepth 1 -type f \
    ! -name 'checksums_${ITER}.md5' \
    -exec md5sum {} \; | sort > '$ITER_TEST_DIR/checksums_${ITER}.md5'"

echo "MD5 checksums saved"

# ============================================
# DATASET SIZE / FILE COUNT
# ============================================

FILE_COUNT=$(sudo find "$ITER_TEST_DIR" -type f | wc -l)
DATASET_SIZE=$(sudo du -sh "$ITER_TEST_DIR" | cut -f1)

echo ""
echo "=========================================="
echo "DATASET ${ITER} CREATED"
echo "=========================================="
echo "Location:     $ITER_TEST_DIR"
echo "Dataset size: $DATASET_SIZE"
echo "File count:   $FILE_COUNT"
echo "=========================================="

# ============================================
# PHASE 4: Delete Test Data
# ============================================

echo ""
echo "=== [ITER $ITER - PHASE 4]: Delete Test Data ==="

BEFORE=$(find "$ITER_TEST_DIR" -type f | wc -l)
echo "Deleting $BEFORE test files..."
sudo rm -rf "$ITER_TEST_DIR"/*
sync
sleep 5
echo "Test data deleted"



# ============================================
# PHASE 5: Measure TRIM Latency
# ============================================

echo ""
echo "=== [ITER $ITER - PHASE 5]: Measure TRIM Latency ==="

    
    # Clear caches
    sync && echo 3 | sudo tee /proc/sys/vm/drop_caches > /dev/null
    
    # Run fstrim with timing
    LOG_FILE="$LOG_DIR/trim_${FILL_LEVEL}_iter${ITER}_$(date +%H%M%S).log"
    /usr/bin/time -v sudo fstrim -v "$MOUNT_POINT" 2>&1 | tee "$LOG_FILE"
    
    # Extract timing
    REAL_TIME=$(grep "Elapsed" "$LOG_FILE" | awk '{print $NF}')
    USER_TIME=$(grep "User time" "$LOG_FILE" | awk '{print $NF}')
    SYS_TIME=$(grep "System time" "$LOG_FILE" | awk '{print $NF}')
    TRIMMED=$(grep -o "[0-9]* bytes" "$LOG_FILE" | tail -1 | sed 's/ bytes//')
    echo "    Real: ${REAL_TIME}s, Trimmed: ${TRIMMED}B"
 
  

# ============================================
# PHASE 6: Attempt Recovery at CURRENT_WAIT
# ============================================

echo ""
echo "=== [ITER $ITER - PHASE 6]: Recovery Attempts ==="
echo "  Waiting ${CURRENT_WAIT} minutes before recovery..."


    # Wait the specified time
    if [ $CURRENT_WAIT -gt 0 ]; then
        echo "    Waiting ${CURRENT_WAIT} minutes..."
        sleep $((CURRENT_WAIT * 60))
    fi
    
    # Try recovery with foremost
    RECOVERY_DIR="$LOG_DIR/recovered_${FILL_LEVEL}_iter${ITER}_${CURRENT_WAIT}min"
    mkdir -p "$RECOVERY_DIR"
    
    # Use foremost to recover deleted files (quick)
    echo "    Running foremost analysis..."
    sudo foremost -i /dev/sda1 -o "$RECOVERY_DIR" -t all 
    
    # Count recovered files (simplified - check if any files appear)
    RECOVERED=0
    if [ -d "$RECOVERY_DIR" ]; then
        RECOVERED=$(find "$RECOVERY_DIR" -type f ! -name "audit.txt" 2>/dev/null | wc -l)
    fi
    
    echo "    Recovered: $RECOVERED files"
    
# Capture timestamp for this run
    RUN_TIMESTAMP=$(date +%Y%m%d_%H%M%S)


echo "$RUN_TIMESTAMP,$ITER,$CURRENT_WAIT,$FILL_LEVEL,$REAL_TIME,$USER_TIME,$SYS_TIME,$TRIMMED,$RECOVERED" >> $LOG_DIR/experiment_results.csv


    echo ""
    echo "Iteration $ITER complete (waited ${CURRENT_WAIT} minutes)"
    echo "---"
    
    # Small delay between iterations
    sleep 2
done

\end{lstlisting}


% %\lstinputlisting[]{} % for source code files directly
% % lstlisting environment for direct inclusion
% \begin{lstlisting}[language=JavaScript]
%     const test = test;
% \end{lstlisting}

% % for computational complexity
% $\mathcal{O}\left(n\log{n}\right)$

% % verbatim
% \verb+numpy+

